% !TEX root = ../main.tex
\begin{ouverture}[derivation]{4}{Variations et optimisation}{Partie 2 · Analyse}
{Signe de la dérivée\enspace·\enspace Tableaux\enspace·\enspace Extremums\enspace·\enspace Modèles}
\objectif{Relier le signe de la dérivée aux variations.}
\objectif{Dresser un tableau avec les valeurs aux bornes.}
\objectif{Déterminer les extremums et comparer deux fonctions.}
\objectif{Résoudre un problème d'optimisation sur un domaine précis.}
\prerequis{Calcul des dérivées, tangentes et second degré.}
\prerequis{Aires, volumes et mise en équation d'un problème.}
\end{ouverture}
\essentiel
\essSection{1}{Du signe de la dérivée au sens de variation}
\begin{theoreme}
Soit $f$ dérivable sur un intervalle $I$.
Si $f'(x)\geq0$ sur $I$, $f$ y est croissante ; si $f'(x)\leq0$, elle y est
décroissante ; si $f'(x)=0$ partout, elle est constante.
Si la dérivée est strictement positive sauf en un nombre fini de points,
la croissance est stricte ; de même pour la décroissance.
\end{theoreme}
\begin{savoirfaire}[Les quatre étapes]
Déterminer le domaine. Calculer et factoriser la dérivée. Étudier son signe.
Tracer les flèches de variation et calculer les valeurs remarquables.
\end{savoirfaire}
\begin{exemple}
Pour $f(x)=x^3-3x$ sur $[-2;2]$, $f'(x)=3(x-1)(x+1)$.
\end{exemple}
\begin{center}\begin{tableau}
\otTab[2.4]{$x$/0.8,$f'(x)$/0.8,$f$/1.8}{$-2$,$-1$,$1$,$2$}
\tkzTabLine{,+,z,-,z,+,}
\tkzTabVar{-/$-2$,+/$2$,-/$-2$,+/$2$}
\end{tableau}\end{center}
\begin{attention}
$f'(a)=0$ ne suffit pas à prouver un extremum : $x^3$ a une dérivée nulle en 0
et reste strictement croissante. Il faut examiner le changement de signe.
Un maximum sur un intervalle peut aussi être atteint à une borne.
\end{attention}
\clearpage
\essSection{2}{Extremums et optimisation}
\begin{propriete}
Lorsque $f'$ passe de $+$ à $-$ en $a$, $f(a)$ est un maximum local.
Lorsqu'elle passe de $-$ à $+$, c'est un minimum local. Pour déterminer un
extremum global sur un segment, comparer également les valeurs aux bornes.
Si $f$ admet un extremum local en un point $a$ intérieur à l'intervalle, alors
$f'(a)=0$ : la tangente en ce point est horizontale.
\end{propriete}
\begin{exemple}
Pour $f(x)=ax^2+bx+c$, $f'(x)=2ax+b$ s'annule en $-b/(2a)$ en changeant de signe.
Si $a>0$, $f$ décroît puis croît : le sommet est un minimum ; si $a<0$, c'est un maximum.
\end{exemple}
\begin{savoirfaire}[Modéliser un problème]
Choisir une variable et son domaine. Exprimer la grandeur à optimiser avec
cette seule variable. Étudier ses variations. Conclure avec les dimensions,
la valeur optimale et l'unité demandées. Pour une quantité entière, comparer
les entiers admissibles voisins de l'optimum réel.
\end{savoirfaire}
\essSection{3}{Comparer deux fonctions et exploiter une symétrie}
\begin{propriete}
Pour comparer $f$ et $g$, étudier le signe de $f-g$. Si $f-g\geq0$, la courbe
de $f$ est au-dessus de celle de $g$.
Sur un domaine symétrique par rapport à 0, $f(-x)=f(x)$ caractérise une fonction
paire ; $f(-x)=-f(x)$, une fonction impaire.
\end{propriete}
\begin{center}
\begin{tikzpicture}\begin{axis}[otLivreGraphe,width=74mm,height=52mm,xmin=-2.2,xmax=2.2,ymin=-0.5,ymax=4,xtick={-2,-1,1,2},ytick={1,2,3}]
\addplot[otNavy,thick,domain=-1.9:1.9]{x^2};
\addplot[only marks,otGreenDark]coordinates{(-1.5,2.25)(1.5,2.25)};
\end{axis}\end{tikzpicture}\hspace{3mm}
\begin{tikzpicture}\begin{axis}[otLivreGraphe,width=74mm,height=52mm,xmin=-2.2,xmax=2.2,ymin=-3,ymax=3,xtick={-2,-1,1,2},ytick={-2,2}]
\addplot[otGreenDark,thick,domain=-2:2]{0.3*x^3};
\addplot[only marks,otNavy]coordinates{(-1.5,-1.0125)(1.5,1.0125)};
\end{axis}\end{tikzpicture}\par
\legende{À gauche : une fonction paire. À droite : une fonction impaire.}
\end{center}
\exercices

\rubrique{Découverte}
\begin{blocLivre}{ex-4-1}
\begin{exercice}[1]{Du signe aux variations}
Une fonction $f$ est dérivable sur $[-3;4]$, avec $f'(x)>0$ sur $[-3;1[$,
$f'(1)=0$ et $f'(x)<0$ sur $]1;4]$.
\begin{questions}\item Donner les variations de $f$ et la position de son maximum.
\item Avec $f(-3)=2$, $f(1)=7$ et $f(4)=-1$, donner aussi son minimum global.
\end{questions}
\end{exercice}
\end{blocLivre}
\begin{blocLivre}{ex-4-2}
\begin{exercice}[1]{Un trinôme sur un segment}
Étudier les variations de $f(x)=x^2-4x+1$ sur $[-1;5]$.
Déterminer son minimum et son maximum sur cet intervalle.
\end{exercice}
\end{blocLivre}
\begin{blocLivre}{ex-4-3}
\begin{exercice}[1]{Une dérivée ne change pas de signe}
Calculer la dérivée de $g(x)=x^3+3x$, justifier ses variations sur $\R$
et déterminer ses extremums sur $[-1;2]$.
\end{exercice}
\end{blocLivre}
\begin{blocLivre}{ex-4-4}
\begin{exercice}[1]{Un quotient simple}
Soit $h(x)=\dfrac{x+1}{x+2}$ sur $[0;4]$.
Calculer $h'$, dresser le tableau de variations et donner les deux extremums.
\end{exercice}
\end{blocLivre}
\begin{blocLivre}{ex-4-5}
\begin{exercice}[1]{Dérivée nulle et extremum}
Pour $f(x)=x^3$, calculer $f'(0)$ et étudier le signe de $f'$.
Le nombre 0 est-il l'abscisse d'un extremum ? Justifier par les variations.
\end{exercice}
\end{blocLivre}
\clearpage
\rubrique{Entraînement}
\begin{blocLivre}{ex-4-6}
\begin{exercice}[2]{Un polynôme du troisième degré}
$f(x)=x^3-3x^2-9x+5$ sur $[-3;5]$.
\begin{questions}\item Montrer que $f'(x)=3(x+1)(x-3)$.
\item Dresser le tableau complet de variations avec les valeurs aux bornes.
\item Distinguer les extremums locaux et les extremums globaux sur $[-3;5]$.
\end{questions}
\end{exercice}
\end{blocLivre}
\begin{blocLivre}{ex-4-7}
\begin{exercice}[2]{Un minimum avec un quotient}
$f(x)=x+4/x$ sur $]0;+\infty[$.
\begin{questions}\item Calculer $f'(x)$ et en étudier le signe.
\item En déduire le minimum de $f$ et l'inégalité $x+4/x\geq4$ pour tout $x>0$.
\item Retrouver cette inégalité en écrivant $f(x)-4$ sous la forme d'un quotient dont le numérateur est un carré.
\end{questions}
\end{exercice}
\end{blocLivre}
\begin{blocLivre}{ex-4-8}
\begin{exercice}[2]{Étudier une fonction paire}
$f(x)=x^4-4x^2$ sur $[-3;3]$.
\begin{questions}\item Vérifier la parité de $f$.
\item Factoriser $f'(x)$ ; étudier les variations sur $[0;3]$.
\item Compléter le tableau sur $[-3;3]$ par symétrie et déterminer les extremums globaux.
\item Vérifier également que $g(x)=x^3-3x$ est impaire sur $\R$ et préciser sa symétrie.
\end{questions}
\end{exercice}
\end{blocLivre}
\begin{blocLivre}{ex-4-9}
\begin{exercice}[2]{Position par rapport à une tangente}
Pour $f(x)=x^3$ sur $[0;3]$, on considère la tangente $T$ en $x=1$.
\begin{questions}\item Écrire l'équation de $T$.
\item Montrer que $x^3-3x+2=(x-1)^2(x+2)$.
\item Comparer la courbe et $T$ sur $[0;3]$. Où se rencontrent-elles ?
\end{questions}
\end{exercice}
\end{blocLivre}
\begin{blocLivre}{ex-4-10}
\begin{exercice}[2]{Une inégalité grâce aux variations}
Soit $f(x)=\sqrt x-\frac x2$ sur $[1;9]$.
\begin{questions}\item Calculer $f'(x)$ pour $x\in]1;9[$ et justifier que $f$ est décroissante.
\item Déterminer ses extremums sur $[1;9]$.
\item En déduire l'encadrement $-\frac32\leq\sqrt x-\frac x2\leq\frac12$ sur cet intervalle.
\end{questions}
\end{exercice}
\end{blocLivre}
\begin{blocLivre}{ex-4-11}
\begin{exercice}[2]{Optimiser une quantité entière}
Pour un nombre entier $n$ d'articles, $0\leq n\leq20$, le bénéfice est
$B(n)=-2n^2+27n-40$ euros.
\begin{questions}\item Étudier d'abord $B(x)$ sur $[0;20]$ pour une variable réelle.
\item Quels sont les deux entiers à comparer pour maximiser $B(n)$ ?
\item Donner la quantité entière optimale et le bénéfice correspondant.
\end{questions}
\end{exercice}
\end{blocLivre}
\clearpage
\rubrique{Problèmes}
\begin{blocLivre}{ex-4-12}
\begin{exercice}[3]{Une boîte sans couvercle}
Dans une feuille carrée de 24 cm de côté, on découpe quatre carrés de côté $x$,
puis on replie les bords. Le volume de la boîte est $V(x)$.
\begin{center}\begin{tikzpicture}[x=1cm,y=1cm]
\fill[otNavyTint](0,0)rectangle(3.7,3.7);
\foreach \a/\b in {0/0,3/0,0/3,3/3}{\fill[white](\a,\b)rectangle++(.7,.7);\draw[otMuted](\a,\b)rectangle++(.7,.7);}
\draw[otNavy,dashed](.7,.7)rectangle(3,3);
\node at(1.85,1.85){fond};\node at(.35,.35){$x$};
\draw[otMuted,<->](0,-.3)--node[below]{$24$ cm}(3.7,-.3);
\end{tikzpicture}\par\legende{Patron schématique, non à l'échelle.}\end{center}
\begin{questions}\item Justifier le domaine $]0;12[$ et établir $V(x)=x(24-2x)^2$.
\item Montrer que $V'(x)=12(x-4)(x-12)$ et étudier les variations.
\item Donner les dimensions de la boîte de volume maximal et ce volume en litres.
\end{questions}
\end{exercice}
\end{blocLivre}
\begin{blocLivre}{ex-4-13}
\begin{exercice}[3]{Un rectangle sous une parabole}
Les sommets supérieurs d'un rectangle sont $(-x;9-x^2)$ et $(x;9-x^2)$,
les deux autres étant sur l'axe des abscisses, avec $0<x<3$.
\begin{questions}\item Faire un schéma et exprimer l'aire $A(x)$.
\item Calculer $A'$ puis déterminer le rectangle d'aire maximale.
\item Donner la largeur, la hauteur et l'aire exactes, puis une valeur approchée de l'aire à $10^{-2}$ près.
\end{questions}
\label{figlivre-r04-13-rectangle-ex}
\figureLivre{r04-13-rectangle}

\end{exercice}
\end{blocLivre}
\begin{blocLivre}{ex-4-14}
\begin{exercice}[3]{Un emballage de périmètre imposé}
Un rectangle a un périmètre de 40 cm. On le fait tourner autour de l'un de
ses côtés pour obtenir un cylindre : le rayon vaut $r$ et la hauteur vaut $h$.
\begin{questions}\item Justifier $h=20-r$ et $0<r<20$.
\item Exprimer le volume $V(r)$ du cylindre et sa dérivée.
\item Déterminer le rayon et la hauteur donnant le volume maximal.
\item Calculer ce volume et expliquer pourquoi les deux limites géométriques
$r=0$ et $r=20$ ne donnent pas un cylindre utilisable.
\end{questions}
\label{figlivre-r04-14-cylindre-ex}
\figureLivre{r04-14-cylindre}

\end{exercice}
\end{blocLivre}
\begin{blocLivre}{ex-4-15}
\begin{exercice}[3]{Distance minimale à un point}
Un point mobile $M(x;x^2)$ se déplace sur la parabole $y=x^2$, avec $x\in[-2;2]$.
On cherche sa distance minimale au point $A(0;2)$.
\partieExo{A}{Éviter de dériver une racine}
\begin{questions}\item Montrer que $AM^2=D(x)=x^4-3x^2+4$.
\item Expliquer pourquoi minimiser $AM$ revient à minimiser $D(x)$.
\end{questions}
\partieExo{B}{Trouver tous les points optimaux}
\begin{questions}[resume]\item Factoriser $D'(x)$ et étudier ses variations sur $[-2;2]$.
\item Déterminer les deux positions de $M$ qui minimisent $AM$.
\item Donner la distance minimale exacte et vérifier qu'elle est inférieure à $AO$.
\end{questions}
\label{figlivre-r04-15-distance-ex}
\figureLivre{r04-15-distance}

\end{exercice}
\end{blocLivre}
